\section{Acquisition consequences and supervised comparison}\label{sec:acquisition-consequences}

The sample requirements in \cref{obj:prop:annotation-threshold} follow from the two terms of the sharp rate. Throughout the acquisition analysis, fix \(d\ge1\), \(\alpha,\beta\in(0,1]\), \(\gamma\in[1,3]\), \(L>1/2\), and an overlap level \(\varepsilon\in(0,1/2)\), as in \cref{obj:thm:sharp-annotation-frontier,obj:thm:supplied-propensity}. Retain the total treatment-record count \(N=n+m\), summed nuisance smoothness \(S=\alpha+\beta\), rate denominator \(\Delta\), smoothness threshold \(S_{\mathrm{crit}}\), and total-record exponent \(q_*\) defined there. The identities below express how labeled responses and treatment records jointly determine precision.

\subsection{Complete acquisition statement}

\begin{propositionv}{prop:annotation-threshold}[Annotation thresholds and acquisition requirements]
Fix public parameters satisfying
\begin{itemize}
\item \textbf{(Dimension.)} \(d\in\mathbb N\) and \(d\ge1\).
\item \textbf{(Nuisance smoothness.)} \(0<\alpha\le1\) and \(0<\beta\le1\).
\item \textbf{(Effect smoothness.)} \(1\le\gamma\le3\).
\item \textbf{(Radius.)} \(L>1/2\).
\item \textbf{(Overlap.)} \(0<\varepsilon<1/2\).
\end{itemize}
Let \(R(n,m)\) and \(R^e(n,m)\) be the minimax and supplied-propensity risks of \cref{obj:def:minimax,obj:def:oracle-risk}, and let \(r_*(n,m)\) be the sharp rate of \cref{obj:def:sharp-rate}, with the fixed parameter arguments suppressed. Define the total treatment-record count \(N\), the summed nuisance smoothness \(S\), the tuning denominator \(\Delta\), the smoothness cutoff \(S_{\mathrm{crit}}\), the total-record exponent \(q_*\), and the oracle rate \(r_o\) by
\[
N=n+m,\qquad S=\alpha+\beta,\qquad
\Delta=2\gamma+d+\frac{\gamma d}{S},
\qquad
S_{\mathrm{crit}}=\frac{\gamma d}{2\gamma+d},
\qquad
q_*=\frac{\gamma d}{S(2\gamma+d)},
\qquad
r_o(n)=n^{-\gamma/(2\gamma+d)}.
\]
All constants below may depend on the fixed public parameters.

For every sequence \(m_n\in\mathbb N\), the following equivalences hold:
\[
\begin{aligned}
\frac{R(n,m_n)}{R^e(n,m_n)}=O(1)
&\quad\Longleftrightarrow\quad
r_*(n,m_n)=O(r_o(n))\\
&\quad\Longleftrightarrow\quad
\exists\,\zeta\in[1,\infty)\\ 
\forall\,n\ \text{ sufficiently large},\\
(n,m_n)\in\mathcal B_*(\zeta)\\
&\quad\Longleftrightarrow\quad
\liminf_{n\to\infty}\frac{n+m_n}{n^{q_*}}>0,
\end{aligned}
\]
where \(\mathcal B_*(\zeta)\) is the region of \cref{obj:def:annotation-region}; the lower limit is understood in the extended nonnegative reals. Moreover,
\[
\begin{aligned}
\frac{R(n,m_n)}{R(n,0)}\longrightarrow0
&\quad\Longleftrightarrow\quad
\frac{r_*(n,m_n)}{r_*(n,0)}\longrightarrow0\\
&\quad\Longleftrightarrow\quad
S<S_{\mathrm{crit}}
\ \text{ and }\\ 
\frac{m_n}{n}\longrightarrow\infty.
\end{aligned}
\]
All sequence limits and order relations are as \(n\to\infty\). If \(S<S_{\mathrm{crit}}\), then \(q_*>1\). For every sequence with \(m_n/n=O(1)\),
\[
R(n,m_n)=O(R(n,0)),
\qquad
R(n,0)=O(R(n,m_n)).
\]
If \(S\ge S_{\mathrm{crit}}\), then every auxiliary-size sequence satisfies
\[
\frac{R(n,m_n)}{R^e(n,m_n)}=O(1).
\]

There exist real constants \(c,C\) with \(0<c\le C\) such that, for every \(n,m\in\mathbb N\) with \(n\ge2\),
\[
c\,r_*(n,m)\le R(n,m)\le C\,r_*(n,m).
\]
For these sample sizes, the exact rate branches are
\[
r_*(n,m)=
\begin{cases}
r_o(n),&S\ge S_{\mathrm{crit}},\\
(nN)^{-\gamma/\Delta},
&S<S_{\mathrm{crit}},\ N\le n^{q_*},\\
r_o(n),
&S<S_{\mathrm{crit}},\ N\ge n^{q_*}.
\end{cases}
\]
At every boundary pair with \(N=n^{q_*}\),
\[
(nN)^{-\gamma/\Delta}=r_o(n).
\]
In particular, for every \(n\ge2\),
\[
r_*(n,0)=
\begin{cases}
n^{-2\gamma/\Delta},&S<S_{\mathrm{crit}},\\
r_o(n),&S\ge S_{\mathrm{crit}}.
\end{cases}
\]
Writing \(s_{\mathrm{pub}}=S/2\) for the average nuisance smoothness, the numerical supervised benchmark coincides with \(r_*(n,0)\): for every \(n\in\mathbb N\),
\[
\begin{cases}
n^{-1/(1+d/(2\gamma)+d/(4s_{\mathrm{pub}}))},
&s_{\mathrm{pub}}<\dfrac{d/4}{1+d/(2\gamma)},\\
n^{-1/(2+d/\gamma)},
&s_{\mathrm{pub}}\ge\dfrac{d/4}{1+d/(2\gamma)}
\end{cases}
=r_*(n,0).
\]
In this identity at \(n=0\), negative powers of zero are assigned the value zero.

Finally, there exist real constants \(c,C>0\), separate from the preceding pair, such that the following acquisition guarantees hold for every target absolute-error scale \(0<\eta<1\) and every \(n,m\in\mathbb N\) with \(n\ge2\). Let \(\mathcal M\) be the primitive class at the fixed parameters from \cref{obj:def:primitive-class}, and use the risk of \cref{obj:def:risk}.

If there exists a Borel decision \(T\) satisfying
\[
\mathcal R_{n,m}(T,P)\le\eta
\qquad\text{for every }P\in\mathcal M,
\]
then
\[
c\,\eta^{-(2+d/\gamma)}\le n,
\qquad
c\,\eta^{-(2+d/\gamma+d/S)}\le n(n+m).
\]
Conversely, if
\[
C\,\eta^{-(2+d/\gamma)}\le n,
\qquad
C\,\eta^{-(2+d/\gamma+d/S)}\le n(n+m),
\]
then the sharp decision \(T_*\) of \cref{obj:def:sharp-decision} is Borel and satisfies
\[
\mathcal R_{n,m}(T_*,P)\le\eta
\qquad\text{for every }P\in\mathcal M.
\]
The exact rate-level acquisition condition is
\[
r_*(n,m)\le\eta
\quad\Longleftrightarrow\quad
\eta^{-(2+d/\gamma)}\le n
\ \text{ and }\
\eta^{-(2+d/\gamma+d/S)}\le n(n+m).
\]
The total treatment-record count also satisfies \(n+m\ge n\).
\end{propositionv}

\subsection{Oracle attainment and strict benefit}

Since \(\Delta=(2\gamma+d)(1+q_*)\), the sharp rate in \cref{obj:def:sharp-rate} satisfies
\[
\frac{r_*(n,m)}{r_o(n)}
=
\max\left\{
1,\,
\left(\frac{n^{q_*}}{N}\right)^{\gamma/\Delta}
\right\}.
\]
Thus, for any finite tolerance \(\zeta\ge1\), the annotation region in \cref{obj:def:annotation-region} has the exact characterization
\[
(n,m)\in\mathcal B_*(\zeta)
\quad\Longleftrightarrow\quad
N\ge \zeta^{-\Delta/\gamma}n^{q_*},
\qquad n\ge2,\quad m\ge0.
\]
The tolerance measures the remaining nuisance contribution relative to the supplied-propensity order. At \(\zeta=1\), the count threshold is the point at which the two terms coincide.

For an auxiliary-size sequence \((m_n)\), write \(N_n=n+m_n\). The uniform risk comparisons in \cref{obj:thm:sharp-annotation-frontier,obj:thm:supplied-propensity} give
\[
\frac{R(n,m_n)}{R^e(n,m_n)}
\asymp
\frac{r_*(n,m_n)}{r_o(n)}.
\]
Consequently, bounded relative risk is equivalent to an eventual positive lower bound on \(N_n/n^{q_*}\), or, equivalently,
\[
\frac{R(n,m_n)}{R^e(n,m_n)}=O(1)
\quad\Longleftrightarrow\quad
\liminf_{n\to\infty}\frac{N_n}{n^{q_*}}>0.
\]
This is the oracle-attainment equivalence in \cref{obj:prop:annotation-threshold}. When \(S\ge S_{\mathrm{crit}}\), the inequality \(q_*\le1\), together with \(N_n\ge n\), places every auxiliary-size sequence in this regime.

Strict benefit compares precision with labeled records alone. In the regime \(S<S_{\mathrm{crit}}\), one has \(q_*>1\) and \(r_*(n,0)=n^{-2\gamma/\Delta}\). Hence
\[
\frac{r_*(n,m_n)}{r_*(n,0)}
=
\max\left\{
n^{-\gamma(q_*-1)/\Delta},\,
\left(\frac{n}{N_n}\right)^{\gamma/\Delta}
\right\}.
\]
The first term tends to zero, and the second tends to zero exactly when \(m_n/n\to+\infty\). For \(S\ge S_{\mathrm{crit}}\), the rate ratio equals one. The uniform minimax comparisons therefore yield the criterion in \cref{obj:prop:annotation-threshold}:
\[
\frac{R(n,m_n)}{R(n,0)}\longrightarrow0
\quad\Longleftrightarrow\quad
S<S_{\mathrm{crit}}
\quad\text{and}\quad
\frac{m_n}{n}\longrightarrow+\infty.
\]
Strict benefit and oracle attainment thus describe distinct acquisition requirements. In the low-smoothness regime, a diverging auxiliary-to-labeled ratio produces a strict improvement, while attainment of the supplied-propensity order requires total treatment information of order at least \(n^{q_*}\). A bounded auxiliary-to-labeled ratio preserves the supervised minimax order.

\subsection{Target-error acquisition}

For a target absolute-error scale \(0<\eta<1\), inversion of the two rate terms gives the exact identity
\[
r_*(n,m)\le\eta
\quad\Longleftrightarrow\quad
\begin{cases}
n\ge \eta^{-(2+d/\gamma)},\\
nN\ge \eta^{-(2+d/\gamma+d/S)}.
\end{cases}
\]
The first inequality specifies the response-label requirement. The second specifies the joint requirement on labeled and treatment-record counts.

To connect these numerical conditions to risk, let \(c_R,C_R>0\) be comparison constants from \cref{obj:thm:sharp-annotation-frontier}, so that
\[
c_Rr_*(n,m)
\le R(n,m)
\le \sup_{P\in\mathcal M}\mathcal R_{n,m}(T_*,P)
\le C_Rr_*(n,m).
\]
A Borel decision whose risk is at most \(\eta\) uniformly over \(\mathcal M\) necessarily satisfies
\[
n\ge (c_R/\eta)^{2+d/\gamma},
\qquad
nN\ge (c_R/\eta)^{2+d/\gamma+d/S}.
\]
Conversely, the sufficient conditions
\[
n\ge (C_R/\eta)^{2+d/\gamma},
\qquad
nN\ge (C_R/\eta)^{2+d/\gamma+d/S}
\]
ensure that the decision \(T_*\) in \cref{obj:def:sharp-decision} has uniform risk at most \(\eta\). Absorbing the fixed powers of these comparison constants gives the necessity and sufficiency constants in \cref{obj:prop:annotation-threshold}.

At the rate level, a fixed labeled count satisfying the first inequality requires
\[
N\ge
\max\left\{
n,\,
\frac{\eta^{-(2+d/\gamma+d/S)}}{n}
\right\}.
\]
In particular, choosing \(n\asymp\eta^{-(2+d/\gamma)}\) gives the total-record requirement
\[
N\gtrsim
\max\left\{
\eta^{-(2+d/\gamma)},\,
\eta^{-d/S}
\right\}.
\]
The smoothness threshold compares these two powers. Below \(S_{\mathrm{crit}}\), the treatment-record requirement grows faster than the labeled requirement; at and above the threshold, the labeled count already supplies the required treatment information.

\subsection{The supervised numerical benchmark}

At \(m=0\), the sharp rate becomes
\[
r_*(n,0)
=
\max\left\{
n^{-\gamma/(2\gamma+d)},\,
n^{-2\gamma/\Delta}
\right\}.
\]
Its exponents and smoothness cutoff coincide with the numerical benchmark in \citet[Section 3, Theorem 1]{KennedyBalakrishnanRobinsWasserman2024}.

The cited theorem concerns a supervised model with a covariate density bounded above by a constant, an \(\alpha\)-smooth propensity, a \(\beta\)-smooth control regression, and a \(\gamma\)-smooth contrast. It gives a minimax mean-absolute-error lower bound for sample sizes larger than a constant depending on \((\alpha,\beta,\gamma,d)\). Its implicit multiplicative constant is positive and independent of sample size, with the density bound and H\"older constants fixed; these are source-model constants. The density restriction accommodates the construction-specific designs used in the source testing argument \citep[Section 3.1]{KennedyBalakrishnanRobinsWasserman2024}. The following identities express the numerical comparison in the notation of this paper.

\begin{lemmav}{lem:published-cate-benchmark}[Published CATE benchmark identities]
Let the dimension \(d\in\mathbb N\) and smoothness orders \(\alpha,\beta,\gamma\in\mathbb R\) satisfy:
\begin{itemize}
\item \textbf{(Dimension.)} \(1\le d\).
\item \textbf{(First nuisance smoothness.)} \(0<\alpha\).
\item \textbf{(Second nuisance smoothness.)} \(0<\beta\).
\item \textbf{(Effect smoothness.)} \(1\le\gamma\).
\end{itemize}
Put \(s_{\mathrm{pub}}=(\alpha+\beta)/2\), the average nuisance smoothness. The numerical exponents and smoothness cutoff displayed in \citet[Section 3, Theorem 1]{KennedyBalakrishnanRobinsWasserman2024} satisfy
\[
\frac{1}{1+d/(2\gamma)+d/(4s_{\mathrm{pub}})}
=\frac{2\gamma}{2\gamma+d+\gamma d/(\alpha+\beta)},
\qquad
\frac{1}{2+d/\gamma}=\frac{\gamma}{2\gamma+d},
\]
and
\[
s_{\mathrm{pub}}<\frac{d/4}{1+d/(2\gamma)}
\quad\Longleftrightarrow\quad
\alpha+\beta<\frac{\gamma d}{2\gamma+d}.
\]
For every \(n\in\mathbb N\) with \(n\ge2\), the numerical piecewise benchmark satisfies
\[
\begin{cases}
n^{-1/(1+d/(2\gamma)+d/(4s_{\mathrm{pub}}))},
& s_{\mathrm{pub}}<(d/4)/(1+d/(2\gamma)),\\
n^{-1/(2+d/\gamma)},
& s_{\mathrm{pub}}\ge(d/4)/(1+d/(2\gamma))
\end{cases}
=
\max\left\{
n^{-\gamma/(2\gamma+d)},
n^{-2\gamma/(2\gamma+d+\gamma d/(\alpha+\beta))}
\right\}.
\]
\end{lemmav}

\Cref{obj:lem:published-cate-benchmark} identifies both the exponents and the boundary between their regimes. Equality at the cutoff places the supervised benchmark at the supplied-propensity order. Under the present public-parameter restrictions, its maximum is exactly \(r_*(n,0)\).

For the known-uniform primitive class in \cref{obj:def:primitive-class}, the minimax interpretation of this expression follows from the original-record lower bounds. The construction in \cref{obj:thm:marked-component-certificate} and the converse in \cref{obj:lem:nuisance-frontier-converse} provide the nuisance contribution in its governing range. The same-class subfamily with \(e_P=\mu_{0,P}=1/2\) in \cref{obj:thm:supplied-propensity} provides the labeled-response floor. Together with \cref{obj:thm:rectangular-upper}, these results establish the sharp risk comparison used throughout the acquisition analysis, while \cref{obj:lem:published-cate-benchmark} supplies the exponent-and-cutoff identity with the published supervised benchmark.
